% Notes and numbers: the figures behind the essays (The Trial).
\newcommand\zcell[1]{\ifnum\lab{sig-#1}=1 \textbf{\lab{z-#1}}\else\lab{z-#1}\fi}
\newcommand\rzcell[1]{\ifnum\lab{rsig-#1}=1 \textbf{\lab{rz-#1}}\else\lab{rz-#1}\fi}
\newcommand\vrow[2]{#2 & \lab{m-#1} & \lab{b-#1} $\pm$ \lab{sd-#1} & \zcell{#1} & \rzcell{#1}\\}
\newcommand\stripz[2]{%
  \begin{tikzpicture}
  \begin{axis}[width=40mm, height=21mm, ymin=0, axis y line=none,
      xmin=\lab{bmin-#1}, xmax=\lab{bmax-#1}, enlarge x limits=0.06, xtick=\empty,
      title={#2}, title style={font=\figfont\scriptsize, yshift=-1.5mm}, clip=false]
    \addplot[const plot mark mid, black, line width=0.4pt] table {../build/the-trial/baseline-#1.dat};
    \draw[line width=1.1pt] (axis cs:\lab{raw-#1},0) -- (axis cs:\lab{raw-#1},\pgfkeysvalueof{/pgfplots/ymax});
    \node[anchor=north, font=\figfont\scriptsize] at (axis cs:\lab{raw-#1},0) {$z=\lab{z-#1}$};
  \end{axis}
  \end{tikzpicture}}

\chapter{Notes and numbers}
\label{ch:notes}

These notes give the figures behind the essays, for readers who want to check
them. Everything here is produced by the program that draws the book, in the
repository that accompanies it.

\section{Scoring}
A sentence's score is the sum of the AFINN-165 scores of its words (F.\,Å.
Nielsen, 2011). Its turn is left for a positive sum, right for a negative one,
and straight for zero. A lexicon word is not scored when it is capitalised and
does not begin a sentence or a quotation, or when the book uses it that way in
at least half of its occurrences (at least three times). For \emph{The Trial}
this excludes \lab{names}.

Sentences are split at \texttt{. ! ?} before a capital, a digit or an opening
quotation mark, and at every paragraph break, but not after initials
(``Josef~K.'') or honorifics (Mr., Mrs.). Chapter headings and subtitles are
not drawn.

\section{Shuffles}
Each order-dependent measure is computed for the book and for \lab{shuffles}
random orderings of the same sentences. $z$ is the book's distance from the
shuffles' mean in standard deviations. The essays call a pattern a property of
the book only when $|z|\ge3$ (bold). With twelve measures per book, this keeps
the chance of a single false alarm per book at about 3\,\%. The last column
repeats every test with \lab{r-words} left unscored.

\begin{table}[h]
\centering\figfont\small
\begin{tabular}{@{}lrrrr@{}}
\toprule
 & \emph{Trial} & shuffled & $z$ & w/o \lab{r-words}\\
\midrule
\vrow{recoil}{Same turn twice (share)}
\vrow{vacillation}{Opposite turns (share)}
\vrow{entropy_h2}{Turn entropy $H_2$ (bits)}
\vrow{corridor_mean}{Mean plain run}
\vrow{corridor_max}{Longest plain run}
\vrow{drift}{Drift}
\vrow{gyration_final}{Spread (gyration)}
\vrow{sprawl}{Sprawl}
\vrow{texture}{Texture (dimension)}
\vrow{crossing_rate}{Crossings / sentence}
\vrow{crossing_gap_median}{Median crossing gap}
\vrow{palimpsest_words}{Overwritten words}
\bottomrule
\end{tabular}
\end{table}

\noindent\stripz{recoil}{Same turn twice}\hfill\stripz{vacillation}{Opposite turns}\hfill\stripz{entropy_h2}{Turn entropy $H_2$}\\[1mm]
\stripz{corridor_mean}{Mean plain run}\hfill\stripz{corridor_max}{Longest plain run}\hfill\stripz{drift}{Drift}\\[1mm]
\stripz{gyration_final}{Spread}\hfill\stripz{sprawl}{Sprawl}\hfill\stripz{texture}{Texture}\\[1mm]
\stripz{crossing_rate}{Crossings / sentence}\hfill\stripz{crossing_gap_median}{Median crossing gap}\hfill\stripz{palimpsest_words}{Overwritten words}

\section{Consecutive sentences}
Pairs of consecutive sentences, by turn of the first (rows) and the second
(columns), against the number expected if order did not matter
(observed\,/\,expected):

\begin{center}\figfont\small
\begin{tabular}{@{}lrrr@{}}
\toprule
 & left & straight & right\\
\midrule
left & \lab{tr-LL} / \lab{te-LL} & \lab{tr-LS} / \lab{te-LS} & \lab{tr-LR} / \lab{te-LR}\\
straight & \lab{tr-SL} / \lab{te-SL} & \lab{tr-SS} / \lab{te-SS} & \lab{tr-SR} / \lab{te-SR}\\
right & \lab{tr-RL} / \lab{te-RL} & \lab{tr-RS} / \lab{te-RS} & \lab{tr-RR} / \lab{te-RR}\\
\bottomrule
\end{tabular}
\end{center}

\section{Spread over time}
\begin{center}
\begin{tikzpicture}
\begin{axis}[width=105mm, height=48mm, xmin=0, xmax=1, ymin=0, xlabel={position in the book}, ylabel={spread (words)},
  label style={font=\figfont\scriptsize}, legend style={draw=none, font=\figfont\scriptsize, at={(0.02,0.98)}, anchor=north west}]
  \addplot[name path=hi, draw=none, forget plot] table[x=t,y=hi] {../build/the-trial/gyration.dat};
  \addplot[name path=lo, draw=none, forget plot] table[x=t,y=lo] {../build/the-trial/gyration.dat};
  \addplot[pattern=north east lines, draw=none] fill between[of=hi and lo];
  \addlegendentry{shuffled, 5--95\,\%}
  \addplot[black, dashed, line width=0.4pt] table[x=t,y=mid] {../build/the-trial/gyration.dat};
  \addlegendentry{shuffled, median}
  \addplot[black, line width=1pt] table[x=t,y=rg] {../build/the-trial/gyration.dat};
  \addlegendentry{\emph{The Trial}}
\end{axis}
\end{tikzpicture}
\end{center}
\noindent The radius of gyration (how far, on average, the path lies from its
centre) as the book is read. It stays within the shuffled band for all but
\lab{rg-above-pct}\,\% of the book.

\section{Definitions}
\begin{description}\small
\item[Knot] the 24\,$\times$\,24-word cell (smoothed over its neighbours) with
the most crossings. A crossing is two strokes that are not neighbours in the
text cutting through each other.
\item[Hub] the sentence crossed by the most others.
\item[Return] the crossing between the two sentences farthest apart in the book.
\item[Coil] the twelve consecutive sentences with the largest net turning.
\item[Corridor] the longest run of plain (unscored) sentences, by words.
\item[Darkest passage] the thirty consecutive sentences with the lowest total score.
\item[Palimpsest] the longest stretch where a later stroke lies exactly on an
earlier one.
\item[Mineshaft] the same walk in three dimensions, tilting \lab{shaft-delta}\textdegree{}
up for each left turn and down for each right turn.
\end{description}
